\chapter{Introduction}

Many questions about real-world phenomena are causal in nature, from medical diagnosis~\cite{richens:2020} to analyzing the effects of policies~\cite{blackwell:2013} or environmental changes~\cite{guevara:2025}. Humans can approach such questions using domain knowledge, experience, and intuition. Answering them automatically, however, requires causal knowledge in a machine-processable form. While causal information can be extracted from large text collections, its unstructured nature makes algorithmic reasoning difficult. Causal graphs provide a structured representation using nodes for concepts or events and directed edges for causal relationships, enabling causal reasoning via graph traversal. The binary question of whether a concept $A$ causes another concept~$B$ can therefore be formulated as a pathfinding problem from $A$ to $B$. Causal graphs such as CauseNet~\cite{heindorf:2020} contain millions of cause--effect relationships extracted from the web, making efficient traversal a challenge.\looseness=-1

A straightforward approach is exhaustive search using Breadth-First Search (BFS), which guarantees completeness and shortest paths in terms of hop count~\cite{cormen:2022}. However, shortest paths are not necessarily the most semantically meaningful. Due to the noisy, automatically extracted nature of large causal graphs, they may include weak or overly general causal links, limiting their usefulness as explanations. For example, in CauseNet, BFS connects \textit{waves} to \textit{destruction} through the path \textit{waves} $\rightarrow$ \textit{death} $\rightarrow$ \textit{destruction}. Although topologically valid, the highly connected intermediate concept \textit{death} provides only a broad connection between the queried concepts. In addition, BFS becomes computationally expensive in graphs with high branching factors, often requiring the exploration of thousands of nodes before reaching the target. To reduce this search space, a reinforcement learning (RL) traversal approach has been proposed~\cite{bluebaum:2024}. This approach learns a policy to guide exploration toward promising regions of the graph, substantially reducing the number of visited nodes. However, maintaining and evaluating a dedicated neural policy introduces architectural complexity and inference overhead.

To make causal question answering more efficient, this thesis replaces the learned traversal policy with a learned search heuristic, removing neural computations from graph traversal at inference time. The proposed approach employs $A^*$ search with the embedding distance between the current and target nodes as its heuristic. A pretrained sentence transformer provides embeddings that guide the search toward semantically closer nodes, reducing unnecessary exploration. Since semantic proximity alone does not necessarily identify the most promising successor, the embeddings are further fine-tuned for the search objective. Training examples are derived from $A^*$ paths that reach the target, teaching the model to favor successors that provide a shorter embedding-space route toward the target. In the previous example, $A^*$ with the fine-tuned embeddings instead finds the equally short path \mbox{\textit{waves} $\rightarrow$ \textit{turbulence} $\rightarrow$ \textit{destruction}}, replacing the broad intermediate concept \textit{death} with the more specific \textit{turbulence}. Empirical evaluation shows that fine-tuning substantially reduces the number of visited nodes while generally retaining competitive effectiveness. BFS remains more effective in some settings, but its search effort increases substantially on larger graphs. Compared with the RL method of Blübaum and Heindorf~\cite{bluebaum:2024}, the proposed approach achieves F\textsubscript{1} scores at least \(7.3\) percentage points higher across all final test settings while being at least \(56\times\) as fast.

The contributions of this thesis are as follows: (1) an embedding-guided \(A^*\) approach is proposed for efficient binary causal question answering over large-scale causal graphs; (2) a search-oriented fine-tuning procedure adapts the embedding space for causal pathfinding, together with Matryoshka Representation Learning to support efficient low-dimensional heuristics; and (3) a systematic evaluation on in-domain and out-of-domain causal graphs and datasets compares the approach with BFS and RL traversal, demonstrating substantial reductions in search effort and runtime while characterizing the resulting effectiveness--efficiency trade-off.

The remainder of this thesis is structured as follows:
The theoretical foundations required for the proposed approach are introduced in Chapter~\ref{chap:foundations}. Related work on causal question answering, causal graphs, and graph traversal methods, including RL-based approaches and learned search heuristics, is reviewed in Chapter~\ref{chap:related_work}. Chapter~\ref{chap:methodology} then presents the problem formulation and the proposed embedding-guided \(A^*\) approach for causal graph traversal. The training procedure, including supervision construction, fine-tuning, validation, and model selection, is described in Chapter~\ref{chap:train}. Experimental evaluation and comparisons with BFS and RL across different datasets and causal graphs follow in Chapter~\ref{chap:results}. Chapter~\ref{chap:discussion} discusses the main findings, trade-offs, and limitations. Finally, Chapter~\ref{chap:conclusion} summarizes the thesis and outlines future work.
